\chapter{Shell and command line}
\label{cap:shell}

\section{Overview}

The \code{src/shell.rs} implements a command-line terminal without dynamic
allocation. It is not a POSIX shell and does not interpret pipes,
redirections, variables, quotes, or scripts. It is a loop that:

\begin{enumerate}
  \item prints the \code{rinux> } prompt;
  \item reads bytes from the PS/2 keyboard or the serial port;
  \item echoes printable characters and handles Backspace;
  \item upon receiving Enter, executes the command;
  \item returns to the beginning.
\end{enumerate}

\section{Main loop}

\begin{lstlisting}[caption={Structure of the shell loop.},label={lst:shellloop}]
pub fn run() -> ! {
    loop {
        crate::print!("rinux> ");
        let mut line = [0u8; LINE_MAX];
        let mut len = 0usize;

        loop {
            let byte = match keyboard::read_char().or_else(serial::read_byte) {
                Some(byte) => byte,
                None => { core::hint::spin_loop(); continue; }
            };

            match byte {
                b'\n' | b'\r' => {
                    println!();
                    execute(&line[..len]);
                    break;
                }
                0x08 | 0x7F => {
                    if len > 0 {
                        len -= 1;
                        crate::print!("\x08");
                    }
                }
                b'\t' => {}
                0x20..=0x7E => {
                    if len < LINE_MAX {
                        line[len] = byte;
                        len += 1;
                        crate::print!("{}", byte as char);
                    }
                }
                _ => {}
            }
        }
    }
}
\end{lstlisting}

\subsection{Line buffer}

The line is a fixed array of 128 bytes on the shell's stack:

\begin{lstlisting}[caption={Line buffer.},label={lst:linebuffer}]
const LINE_MAX: usize = 128;
let mut line = [0u8; LINE_MAX];
let mut len = 0usize;
\end{lstlisting}

When the line is full, new characters are silently ignored. There is no
beep, error message, or dynamic growth.

\subsection{Input and echo}

Input comes from two sources, in this order:

\begin{enumerate}
  \item \code{keyboard::read\_char()}: PS/2 keyboard in the VGA window;
  \item \code{serial::read\_byte()}: serial terminal, when available.
\end{enumerate}

If both return \code{None}, the shell executes \code{spin\_loop()} and tries
again. There is no \emph{sleep}, since no timer exists.

Printable characters (0x20 to 0x7E) are stored in the buffer and echoed to
the console. Backspace (\code{0x08} or \code{0x7F}) decrements the length and
prints \code{\textbackslash x08}, which the VGA console interprets as erasing
the previous cell. Tab (\code{\textbackslash t}) is ignored.

\section{Parsing}

\subsection{Trim}

Before comparing, the line goes through \code{trim}, which removes ASCII
whitespace at the beginning and end. The function returns a \emph{slice} of
the line itself, without copying:

\begin{lstlisting}[caption={Removal of whitespace at the edges.},label={lst:trim}]
fn trim(bytes: &[u8]) -> &[u8] {
    let start = bytes.iter()
        .position(|byte| !byte.is_ascii_whitespace())
        .unwrap_or(bytes.len());
    let end = bytes.iter()
        .rposition(|byte| !byte.is_ascii_whitespace())
        .map_or(start, |index| index + 1);
    &bytes[start..end]
}
\end{lstlisting}

An empty line or one consisting only of spaces executes nothing.

\subsection{Text conversion}

\code{as\_text} tries to interpret the bytes as UTF-8. Since the console only
accepts ASCII, in practice the conversion always succeeds for typed text; if
it fails, it returns the fixed string \code{"<bytes invalidos>"}.

\subsection{Commands}

Execution is a chain of \code{if}/\code{else if} comparisons. There is no
dispatch table, argument \emph{parsing}, or syntax analysis.

\begin{table}[H]
\centering
\caption{Implemented commands.}
\label{tab:comandos}
\begin{tabularx}{\linewidth}{@{}lX@{}}
\toprule
\textbf{Command} & \textbf{Effect} \\
\midrule
\code{help} & Lists the available commands. \\
\code{clear} & Clears the VGA screen and repositions the cursor. \\
\code{echo <texto>} & Prints the text after the first space. \\
\code{info} & Shows kernel and input model information. \\
\code{reboot} & Restarts the PC via command \code{0xFE} to the 8042. \\
\code{halt} & Prints a message and stops the CPU with \code{cli; hlt}. \\
\bottomrule
\end{tabularx}
\end{table}

\begin{nota}[Behavior details]
\begin{itemize}
  \item \code{echo} without an argument is not recognized as a command; only
        \code{echo } followed by text matches the prefix.
  \item \code{echo   ola} preserves the internal spaces after the
        \code{"echo "} prefix; \emph{trim} removes only the edges.
  \item Comparison is case-sensitive: \code{HELP} does not work.
  \item Unknown commands produce two lines of message.
\end{itemize}
\end{nota}

\section{\code{info} output}

The \code{info} command prints:

\begin{lstlisting}[style=shell,caption={Output of the \code{info} command.},label={lst:info}]
Kernel: Rust, no_std, x86_64
Boot:   Multiboot2 (BIOS) -> modo longo
Video:  VGA texto 80x25, buffer 0xB8000
Entrada: teclado PS/2 e serial COM1 (polling, sem IRQ)
Heap:   nenhum (apenas core, sem alocacao)
\end{lstlisting}

This information is static; there is no reading of CPU, memory, or time.

\section{Special commands}

\subsection{\code{halt}}

Prints \code{"CPU parada. Use reset ou feche o QEMU."} and calls
\code{arch::halt\_loop()}, which executes \code{cli; hlt} indefinitely. The
only way out is a physical reset or closing the emulator.

\subsection{\code{reboot}}

Prints \code{"Reiniciando..."} and calls \code{arch::reboot()}. The reset may
not work in all environments; if the command to the 8042 is ignored, the CPU
ends up in \code{halt\_loop()}.

\section{How to extend}

To add a command:

\begin{enumerate}
  \item add an \code{else if line == b"meucomando"} branch in
        \code{execute};
  \item use \code{println!} for output;
  \item if you need arguments, use \code{strip\_prefix} as in
        \code{echo};
  \item keep the command allocation-free and non-blocking;
  \item update the \code{help} list and this documentation.
\end{enumerate}

Example:

\begin{lstlisting}[caption={Skeleton of a new command.},label={lst:novocomando}]
} else if line == b"versao" {
    println!("Rust VGA Kernel 0.1.0");
} else if let Some(arg) = line.strip_prefix(b"upper ") {
    for byte in arg {
        crate::print!("{}", byte.to_ascii_uppercase() as char);
    }
    println!();
}
\end{lstlisting}

\section{Limitations}

\begin{itemize}
  \item No history, arrows, Ctrl+C, Ctrl+D, or editing in the middle of the
        line.
  \item No quotes, escaping, pipes, redirection, or variables.
  \item No structured arguments: each command decides how to interpret the
        remainder.
  \item No program execution: there is no \emph{loader} or processes.
  \item No UTF-8 translation: typed accents may become \code{?}.
  \item The line is limited to 128 bytes and excess is silently discarded.
\end{itemize}
